\section*{अध्याय \ref{ch:b1:observing-the-sky} --- दूरबीनें, संसूचक और प्रकाशमिति}

\begin{solution}{exo:b1:observing-the-sky:1}
5 कांतियों का अंतर \omterm{def:b1:observing-the-sky:flux}{फ़्लक्स} में 100 गुने का अंतर है: कांति 1.0 वाला तारा 100 गुना अधिक \omterm{def:b1:observing-the-sky:flux}{फ़्लक्स} भेजता है।
रीजेल और पोलारिस: $10^{0.4\times(2.02 - 0.13)} =
10^{0.756} = 5.70$; रीजेल, जिसकी कांति छोटी है, अधिक चमकीला है।
\end{solution}

\begin{solution}{exo:b1:observing-the-sky:2}
$d = 1/0.19495 = \qty{5.13}{\parsec}$, अर्थात् $5.1295\times206265 =
\qty{1.058e6}{\au}$ और $5.1295\times3.2616 = \qty{16.7}{\lightyear}$।
\end{solution}

\begin{solution}{exo:b1:observing-the-sky:3}
$1.22\times\num{550e-9}/2.4 = \qty{2.80e-7}{rad} = \ang{;;0.058}$, धुंधलेपन \ang{;;0.40} से लगभग
सात गुना बारीक: अंतरिक्ष में दूरबीन अपनी विवर्तन-सीमा तक पहुँचती है। \qty{6.5}{m} के लिए
\qty{2.0}{\micro\metre} पर: $1.22\times\num{2.0e-6}/6.5 = \qty{3.75e-7}{rad} = \ang{;;0.077}$।
\end{solution}

\begin{solution}{exo:b1:observing-the-sky:4}
$B - V$ रीजेल ($-0.03$) से अल्टेयर (0.22), पोलारिस (0.60) और बेटेल्जूज़ (1.85) तक बढ़ता है: सबसे गर्म
रीजेल है, फिर अल्टेयर, पोलारिस और बेटेल्जूज़। दूरी दोनों बैंडों को एक ही गुणक से मंद करती है, इसलिए
$B - V$ उससे अप्रभावित रहता है (जब तक रास्ते में धूल प्रकाश को लाल न करे)।
\end{solution}

\begin{solution}{exo:b1:observing-the-sky:5}
$X_1 = \sec\ang{30} = 1.1547$, $X_2 = \sec\ang{60} = 2.0000$।
$k = (6.590 - 6.412)/(2.0000 - 1.1547) = 0.178/0.8453 = \num{0.211}$ प्रति \omterm{def:b1:observing-the-sky:extinction}{वायु-राशि}, और
$m_0 = 6.412 - 0.2106\times1.1547 = 6.169$।
\end{solution}

\begin{solution}{exo:b1:observing-the-sky:6}
$\mathrm{SNR} = \sqrt{400} = 20$। SNR $\sqrt t$ के अनुपात में बढ़ता है: 100 के लिए,
$t = \qty{10}{s}\times(100/20)^2 = \qty{250}{s}$। \num{0.01} कांति की त्रुटि के लिए
$\mathrm{SNR} = 1.086/0.01 = 108.6$ चाहिए, इसलिए
$t = 10\times(108.6/20)^2 = \qty{295}{s}$।
\end{solution}

\begin{solution}{exo:b1:observing-the-sky:7}
दूरी $d = 1/0.00378 = \qty{264.6}{\parsec}$ है, और \omterm{def:b1:observing-the-sky:absolute}{निरपेक्ष कांति}
$M = 0.13 - 5\log_{10}26.46 = 0.13 - 7.11 = -6.98$। $L/L_\odot = 10^{0.4\times(4.83 + 6.98)} =
10^{4.724} = \num{5.3e4}$ V में।
\end{solution}

\begin{solution}{exo:b1:observing-the-sky:8}
$\pi\times1.2^2 = \qty{4.52}{m^2}$, और $25/4.52 = 5.5$: वेब दूरबीन उतने ही फ़ोटॉन 5.5 गुना छोटे
अनावरण में एकत्र करती है। \qty{6.5}{m} आर-पार की एक पूरी चकती का $\pi\times3.25^2 = \qty{33.2}{m^2}$
होगा; वास्तविक दर्पण उसका $25/33.2 = \qty{75}{\percent}$ है, क्योंकि 18 षट्कोण उस चकती के किनारे को
अनाच्छादित छोड़ते हैं, और खंडों के बीच के अंतराल तथा केंद्रीय छिद्र कुछ एकत्र नहीं करते।
\end{solution}

\begin{solution}{exo:b1:observing-the-sky:9}
(क) $\qty{10}{\parsec} = \qty{3.0857e17}{m}$, और इस प्रकार $F = \num{3.828e26}/(4\pi
\times(\num{3.0857e17})^2) = \num{3.828e26}/\num{1.1965e36} =
\qty{3.20e-10}{W.m^{-2}}$, $1361/\num{3.20e-10} = \num{4.3e12}$ गुना कम, पृथ्वी की तुलना में।
(ख) $5\log_{10}(d/\qty{10}{\parsec}) = -26.74 - 4.83 = -31.57$, अतः
$d = \qty{10}{\parsec}\times10^{-6.314} = \qty{4.853e-6}{\parsec}$, अर्थात्
$\num{4.853e-6}\times206265 = \qty{1.001}{\au}$।
\end{solution}

\begin{solution}{exo:b1:observing-the-sky:10}
वायुमंडल के ऊपर \omterm{def:b1:observing-the-sky:flux}{फ़्लक्स} $F_0$ है; नीचे, $F = F_0e^{-\tau_0X}$। तब
$m - m_0 = -2.5\log_{10}(F/F_0) = 2.5\log_{10}e\;\tau_0X$, जो $kX$ है, जिसमें
$k = 2.5\log_{10}e\;\tau_0 = 1.0857\,\tau_0$। $k = 0.160$ के लिए:
$\tau_0 = 0.160/1.0857 = 0.147$। \omterm{def:b1:spherical-astronomy:celestial-sphere}{खमध्य} पर प्रकाश का $e^{-0.147} = 0.863$ पहुँचता है, \qty{86}{\percent};
$X = 2.31$ पर, $e^{-0.147\times2.31} =
e^{-0.340} = 0.712$, \qty{71}{\percent}। $z = \ang{85}$ पर $\sec z =
11.5$; पर समतल-समांतर प्रतिरूप क्षितिज पर पथ को अनंत बना देता है, जबकि वक्र वायुमंडल से
गुज़रने वाला पथ परिमित रहता है, और अपवर्तन किरण को भी मोड़ता है: क्षितिज के निकट सूत्र \omterm{def:b1:observing-the-sky:extinction}{वायु-राशि}
को बढ़ाकर बताता है और उसकी त्रुटियाँ तेज़ी से बढ़ती हैं।
\end{solution}

\begin{solution}{exo:b1:observing-the-sky:11}
(क) $\mathrm{SNR} = Rt/\sqrt{(R + B)t}$, अतः $t = \mathrm{SNR}^2(R + B)/R^2
= 100\times2050/2500 = \qty{82}{s}$। (ख) \omterm{def:b1:observing-the-sky:seeing}{धुंधलापन} आधा करने पर प्रतिबिंब का क्षेत्रफल, और इसलिए
$B$, 4 से भाग खाता है: $B = 500$ और $t = 100\times550/2500 =
\qty{22}{s}$। (ग) $B \gg R$ के लिए, $t \approx \mathrm{SNR}^2B/R^2$। स्रोत की दर $R$, $D^2$ के
साथ बढ़ती है; आकाश की दर $B$ संग्रहण क्षेत्रफल $D^2$ और प्रतिबिंब के आकाश-क्षेत्रफल $\theta^2$ के
गुणनफल के साथ बढ़ती है। अतः $t \propto
\theta^2D^2/D^4 = \theta^2/D^2$: \omterm{def:b1:observing-the-sky:seeing}{धुंधलापन} आधा करना व्यास दुगना करने जितना मूल्यवान है।
\end{solution}

\begin{solution}{exo:b1:observing-the-sky:12}
(क) $N = 1200\times100 = \num{120000}$ रेखाएँ; पहली कोटि में विभेदन-क्षमता \num{1.2e5} है, और
$\Delta\lambda = 550/(\num{1.2e5}) =
\qty{4.6e-3}{nm}$। (ख) $v = c\,\Delta\lambda/\lambda = c/(\num{1.2e5}) =
\qty{2.5}{km.s^{-1}}$। (ग) $\Delta\lambda = 550\times29.78/\num{299792} =
\qty{0.055}{nm}$, लगभग 12 विभेदन तत्व। प्रेक्षक को प्रेक्षण की तिथि और समय के लिए, तारे की ओर पृथ्वी
के कक्षीय वेग का घटक परिकलित करके उसे घटाना चाहिए, ताकि तारे का सूर्य के सापेक्ष वेग मिले।
\end{solution}

\begin{solution}{pb:b1:observing-the-sky:1}
\textbf{1.} $d = 1/0.37921 = \qty{2.637}{\parsec}$।

\textbf{2.} $2.637\times206265 = \qty{5.44e5}{\au}$;
$2.637\times\num{3.0857e16} = \qty{8.14e16}{m}$।

\textbf{3.} $2.637\times3.2616 = \qty{8.60}{\lightyear}$: आज रात देखा गया प्रकाश सिरियस से 8.6 वर्ष पहले चला था।

\textbf{4.} दिशा $2p = \ang{;;0.758}$ से झूलती है, जो धुंधलेपन का लगभग दोगुना है। झूलना प्रतिबिंब के
केंद्र की स्थिति से मापा जाता है, उसके किनारे से नहीं: बहुत-से फ़ोटॉनों पर औसत करने पर केंद्र
प्रतिबिंब की चौड़ाई के एक छोटे भाग तक पाया जाता है, और जितने अधिक फ़ोटॉन गिने जाएँ उतनी ही
सूक्ष्मता से।

\textbf{5.} $\Delta d/d = \Delta p/p = 0.001/0.37921 = \qty{0.26}{\percent}$,
अर्थात् $8.60\times0.0026 = \qty{0.023}{\lightyear}$।

\textbf{6.} $10^{0.4\times(0.03 + 1.46)} = 10^{0.596} = 3.94$।

\textbf{7.} $10^{0.4\times(26.74 - 1.46)} = 10^{10.112} = \num{1.29e10}$।

\textbf{8.} $10^{0.4\times(8.44 + 1.46)} = 10^{3.96} = \num{9120}$।

\textbf{9.} $z = \ang{90} - \ang{82.09} = \ang{7.91}$, $X = \sec\ang{7.91} =
1.0096$; क्षीणन $0.15\times1.0096 = 0.151$ कांति; मापी गई
$m = -1.46 + 0.151 = -1.31$।

\textbf{10.} $z = \ang{60}$, $X = 2.000$; क्षीणन 0.30 कांति,
$m = -1.16$। संचरित भाग $10^{-0.4\times0.30} = 0.759$ है:
\qty{24}{\percent} प्रकाश खो जाता है।

\textbf{11.} उदय पर $z = \ang{90}$ और $\sec z$ अनंत है, जबकि वक्र वायुमंडल से गुज़रता वास्तविक पथ लंबा
परंतु परिमित है; अपवर्तन भी किरण को मोड़ता है। क्षितिज के निकट क्षीणन बड़ा है और तेज़ी से बदलता है,
और समतल-समांतर सुधार गलत है।

\textbf{12.} $\num{3.0e4}\times10^{0.4\times(10.0 - 8.44)} = \num{3.0e4}
\times10^{0.624} = \num{3.0e4}\times4.21 = \num{1.26e5}$ प्रकाश-इलेक्ट्रॉन प्रति सेकंड।

\textbf{13.} 1000 का SNR $N = 1000^2 = \num{1.0e6}$ प्रकाश-इलेक्ट्रॉन चाहता है:
$t = \num{1.0e6}/\num{1.26e5} = \qty{7.9}{s}$। कांति की त्रुटि $1.086/1000 = \num{0.0011}$ है।

\textbf{14.} $\num{3.0e4}\times10^{0.4\times(10.0 + 1.46)} = \num{3.0e4}
\times10^{4.584} = \num{1.15e9}$ प्रकाश-इलेक्ट्रॉन प्रति सेकंड।

\textbf{15.} $0.20\times\num{1.15e9} = \num{2.30e8}$ प्रति सेकंड सबसे चमकीले पिक्सेल में, जो
$\num{1.0e5}/\num{2.30e8} =
\qty{4.3e-4}{s}$ में भर जाता है, अर्धमिलीसेकंड से कम।

\textbf{16.} सिरियस के पंख $\num{1.15e9}\times10^{-4} =
\num{1.15e5}$ प्रति सेकंड देते हैं; सिरियस~B $0.20\times\num{1.26e5} =
\num{2.52e4}$ प्रति सेकंड। सिरियस की चकाचौंध सिरियस~B के शिखर से 4.6 गुना उज्ज्वल है।

\textbf{17.} पिक्सेल प्रति सेकंड $\num{1.15e5} + \num{2.52e4} = \num{1.40e5}$ प्राप्त करता है और
$\num{1.0e5}/\num{1.40e5} = \qty{0.71}{s}$ में भर जाता है। \qty{0.5}{s} में: सिग्नल
$\num{2.52e4}\times0.5 = \num{1.26e4}$, कुल
$\num{7.01e4}$, शोर $\sqrt{\num{7.01e4}} = 265$, इसलिए $\mathrm{SNR} =
\num{1.26e4}/265 = 48$। फ़ोटॉन शोर की दृष्टि से सिरियस~B पहचाना जा सकता है; उसे कठिन यह बनाता है
कि चकाचौंध का मान उसके पाँचवें भाग से कहीं बेहतर सटीकता से जानकर घटाना होगा, जबकि सिरियस का केंद्र
एक मिलीसेकंड से कम में संतृप्त हो जाता है।

\textbf{18.} $m - M = 5\log_{10}(2.637/10) = 5\times(-0.579) = -2.89$।

\textbf{19.} $M = -1.46 + 2.89 = 1.43$।

\textbf{20.} $d = 1/0.13023 = \qty{7.679}{\parsec}$;
$m - M = 5\log_{10}0.7679 = -0.57$, अतः $M = 0.03 + 0.57 = 0.60$।

\textbf{21.} वेगा, जिसकी \omterm{def:b1:observing-the-sky:absolute}{निरपेक्ष कांति} छोटी है, गुणक
$10^{0.4\times(1.43 - 0.60)} = 10^{0.332} = 2.15$ से। सिरियस केवल इसलिए अधिक चमकीला दिखता है क्योंकि वह निकट है।

\textbf{22.} $m = M_\odot + (m - M) = 4.83 - 2.89 = 1.94$: सूर्य एक काफ़ी चमकीला नंगी-आँख का तारा होगा।

\textbf{23.} $5\log_{10}(d/\qty{10}{\parsec}) = 6.0 - 1.43 = 4.57$, अतः
$d = 10\times10^{0.914} = \qty{82.0}{\parsec}$, लगभग \qty{267}{\lightyear}।

\textbf{24.} $L/L_\odot = 10^{0.4\times(4.83 - 1.43)} = 10^{1.36} = 22.9$:
V बैंड में सिरियस \textbf{लगभग 23 सूर्यों} से अधिक चमकता है।
\end{solution}
